%==============================================================================
%  config/commands.tex  --  Colores, entornos y macros propias
%  Tesis doctoral - Universidad Nacional de San Luis (UNSL)
%==============================================================================

% ---- Colores ----------------------------------------------------------------
\definecolor{codebg}{rgb}{0.97,0.97,0.97}   % fondo de los bloques de codigo
\definecolor{codeframe}{rgb}{0.80,0.80,0.80}

% ---- Configuracion de minted (codigo fuente) --------------------------------
% Se define aca porque depende de los colores de arriba. minted ya fue cargado
% en config/packages.tex.
\setminted{
  fontsize=\small,
  breaklines=true,        % corta lineas largas
  breakanywhere=true,
  tabsize=4,
  frame=lines,
  framesep=2mm,
  bgcolor=codebg,
  linenos=false,
}
\newcommand{\code}[1]{\mintinline{text}{#1}}   % codigo corto en linea

% ---- Referencias cruzadas en espanol (cleveref) -----------------------------
\crefname{figure}{figura}{figuras}
\Crefname{figure}{Figura}{Figuras}
\crefname{table}{tabla}{tablas}
\Crefname{table}{Tabla}{Tablas}
\crefname{equation}{ecuacion}{ecuaciones}
\Crefname{equation}{Ecuacion}{Ecuaciones}
\crefname{chapter}{capitulo}{capitulos}
\Crefname{chapter}{Capitulo}{Capitulos}
\crefname{section}{seccion}{secciones}
\Crefname{section}{Seccion}{Secciones}
\crefname{algorithm}{algoritmo}{algoritmos}
\Crefname{algorithm}{Algoritmo}{Algoritmos}

% ---- Operadores y delimitadores matematicos ---------------------------------
\DeclareMathOperator*{\argmax}{arg\,max}
\DeclareMathOperator*{\argmin}{arg\,min}
\DeclarePairedDelimiter{\ceil}{\lceil}{\rceil}
\DeclarePairedDelimiter{\floor}{\lfloor}{\rfloor}
\DeclarePairedDelimiter{\abs}{\lvert}{\rvert}

% ---- Entornos de teoremas (numerados por capitulo, en espanol) --------------
\theoremstyle{plain}
\newtheorem{teorema}{Teorema}[chapter]
\newtheorem{proposicion}[teorema]{Proposicion}
\newtheorem{corolario}[teorema]{Corolario}
\newtheorem{lema}[teorema]{Lema}
\newtheorem{conjetura}[teorema]{Conjetura}

\theoremstyle{definition}
\newtheorem{definicion}[teorema]{Definicion}
\newtheorem{ejemplo}[teorema]{Ejemplo}

\theoremstyle{remark}
\newtheorem{observacion}[teorema]{Observacion}
\newtheorem{nota}[teorema]{Nota}

\renewcommand{\proofname}{Demostraci\'on}   % entorno proof en espanol

% rac.sty pone el numero de capitulo en romano; forzamos numeracion arabe
% para que los teoremas queden "1.2" en lugar de "I.2" (consistente con secciones).
\renewcommand{\theteorema}{\arabic{chapter}.\arabic{teorema}}

% ---- Palabras clave de algoritmos en espanol (solo si el paquete esta) ------
\makeatletter
\@ifpackageloaded{algpseudocode}{%
  \floatname{algorithm}{Algoritmo}%
  \renewcommand{\algorithmicrequire}{\textbf{Entrada:}}%
  \renewcommand{\algorithmicensure}{\textbf{Salida:}}%
  \renewcommand{\algorithmicif}{\textbf{si}}%
  \renewcommand{\algorithmicthen}{\textbf{entonces}}%
  \renewcommand{\algorithmicelse}{\textbf{si no}}%
  \renewcommand{\algorithmicend}{\textbf{fin}}%
  \renewcommand{\algorithmicwhile}{\textbf{mientras}}%
  \renewcommand{\algorithmicdo}{\textbf{hacer}}%
  \renewcommand{\algorithmicfor}{\textbf{para cada}}%
  \renewcommand{\algorithmicforall}{\textbf{para todo}}%
  \renewcommand{\algorithmicreturn}{\textbf{retornar}}%
  \renewcommand{\algorithmicprocedure}{\textbf{procedimiento}}%
  \renewcommand{\algorithmicfunction}{\textbf{funcion}}%
  \algtext*{EndWhile}% quita el "fin mientras"
  \algtext*{EndIf}%    quita el "fin si"
  \algtext*{EndFor}%   quita el "fin para"
}{}
\makeatother

% ---- Marca de pendiente para borradores -------------------------------------
% Uso:  \todo{revisar esta cifra}
\newcommand{\todo}[1]{%
  \textbf{\textcolor{red}{[TODO: #1]}}%
  \marginpar{\textcolor{red}{\tiny$\blacktriangleleft$ TODO}}%
}

% ---- Entorno de dedicatoria -------------------------------------------------
\newenvironment{dedicatoria}
  {\clearpage \thispagestyle{empty} \vspace*{\stretch{1}} \slshape \raggedleft}
  {\par \vspace{\stretch{3}} \clearpage}

% ---- Encabezado para secciones preliminares sin número (Resumen, Abstract) --
% Replica el estilo centrado de rac.sty y lo agrega al índice.
\newcommand{\frontsection}[1]{%
  \newpage
  \addcontentsline{toc}{chapter}{#1}%
  \hbox{}\twoinmar
  \begin{center}\large\bfseries #1\end{center}%
  \vspace{0.4in}%
}
